\chapter{Gröbner Bases}

This chapter provides a concise introduction to Gröbner bases and Buchberger's algorithm. The discussion is tailored for an audience of beginning graduate students or advanced undergraduates in mathematics and computer science. Accordingly, we assume a familiarity with standard undergraduate curricula, like calculus and linear algebra, but while a background in abstract algebra, commutative algebra, or algebraic geometry is advantageous, it is not a prerequisite for following the exposition provided here. No original work presented.

For those seeking a more comprehensive treatment, the first two chapters of \cite{cox2015_IVA} offer an excellent and widely recognized alternative introduction. Additional foundational perspectives can be found in \cite{cox2005_UAG, Kreuzer2000_CAG, Sturmfels2005_WhatIsGB}, while \cite{Adams1994_IGB} provides a more advanced theoretical framework. From a historical perspective, this work draws upon Bruno Buchberger's seminal doctoral thesis \cite{Buchberger1965_algorithmGerman, Buchberger1965_algorithmEnglish} and his subsequent survey \cite{Buchberger1985_SurveyArticle}. We also acknowledge modern optimizations that have fundamentally shaped the efficiency of the algorithm, such as the work of \cite{Gebauer1988_InstallationOfGB, Giovini1991_OneSugarCube, Sun2011_F5InBuchbergersStyle}, which are well contextualized among other relevant works in the introductory material of \cite{Eder2012_SignatureBasedAlgorithms}.

The primary objective of this chapter is to establish an understanding of Buchberger's algorithm, with a particular emphasis on the role that $S$-pair selection strategies play in its computational performance. As the ultimate goal of this thesis is to leverage machine learning to optimize these selection processes, establishing this mathematical foundation is a critical first step.